\section{Introducción}
\label{sec:intro}

\IEEEpaperstring{La suplantación de identidad y la presentación de documentos falsificados} (tales como fotocopias a color, capturas en pantallas o impresiones de alta fidelidad) constituyen uno de los mayores vectores de fraude en las industrias financiera y de tecnología. El presente software, denominado **NeoAstrum**, propone un enfoque robusto que combina la agilidad y accesibilidad de la aplicación de mensajería **WhatsApp** como interfaz conversacional, con el poder analítico de la computación en la nube.

\subsection{Alcance del Sistema}
El sistema guía al usuario de manera intuitiva para que proporcione secuencialmente las fotos de la parte frontal y trasera de su cédula colombiana. Una vez capturadas, el motor de orquestación coordina la descarga de las imágenes desde los servidores de Meta, las transfiere a un microservicio en AWS Lambda y realiza la validación técnica. 

\subsection{Objetivos del Software}
\begin{itemize}
    \item **Accesibilidad:** Permitir a cualquier ciudadano validar su identidad a través de WhatsApp sin necesidad de instalar aplicaciones adicionales.
    \item **Automatización:** Retornar una clasificación de veracidad (Real vs Fraude/Fotocopia) en menos de 5 segundos.
    \item **Telemetría y Monitoreo:** Registrar cada etapa del flujo con tiempos de ejecución e indicadores clave para su posterior explotación en un panel web interactivo.
\end{itemize}
